\def\CLASSINPUTtoptextmargin{0.78in}
\def\CLASSINPUTbottomtextmargin{1.06in}
\documentclass[conference]{IEEEtran}
\usepackage{amsmath,amssymb,graphicx,url}
\usepackage{booktabs}
\newcommand{\indicator}{\mathbf{1}}
\begin{document}
\title{[Experiment, Analysis, and Benchmark]\\Auditing LLM-Based Semantic Aggregation\\over Fixed Candidate Pools}
\author{\IEEEauthorblockN{Author information to be completed before submission}\IEEEauthorblockA{Names, affiliations, email addresses, and ORCIDs pending}}
\maketitle
\begin{abstract}
Language-model predicates turn retrieved records into semantic aggregates, but completion, label fidelity, ratio accuracy, and execution cost measure different properties. We audit these properties on fixed candidate pools using pointwise, batched, specification-based, and disagreement-triggered plans. The study covers two product-relevance datasets and adapted biomedical and legal workloads with three models, retaining failed outputs and simple, supervised, embedding, and LOTUS references. @@CONTROLLED_ABSTRACT@@ Post-hoc analysis of 302 product pools shows that an apparent routing benefit need not justify additional judgments. On the 250-pool ESCI selective sample, Qwen's checker has ratio MAE 0.373 versus 0.338 for an OR combination of the same two keyed outputs, which avoids 739 pointwise calls; an all-positive predictor attains 0.185. SmolLM2's selective outputs equal an AND rule on both product samples because every disagreement check returns negative. At equal pointwise-label budgets, Qwen's improvement over uniform checking is concentrated in pools with only one reference class. We provide versioned inputs, executable contracts, recorded traces, and independently checked analyses. The resulting audit separates structural validity from task accuracy and tests whether checking earns its cost. Its conclusions concern fixed relevance proxies and the evaluated configurations, rather than user satisfaction or deployment-scale superiority.

\end{abstract}
\begin{IEEEkeywords}
semantic query processing, semantic aggregation, execution contracts, structured output, data quality, large language models
\end{IEEEkeywords}

\section{Introduction}
A semantic query may ask what fraction of retrieved products meets a request. Its denominator is the candidate list; its predicate may be a language model. If six candidates receive Y and two receive no valid judgment, averaging the returned labels produces 1, whereas the fixed-pool ratio can be as low as $6/8$ even if the six accepted labels are correct. Requiring all eight outputs repairs this accounting error. It does not show that the labels, or their aggregate, are accurate.

Semantic engines already optimize model calls, batch records, and check uncertain decisions \cite{lotus,docetl,palimpzest,sema,nousdb}. Their evaluation must distinguish the property each check establishes. A structurally valid output can attach a wrong label to a valid ID. A repeated decision can agree and still be wrong. Conversely, false positives and false negatives can cancel: an exact ratio succeeds at a ratio-only task even when its labels are unsuitable for reuse. Finally, a checking policy can beat a random allocation of calls while losing to a rule that makes no additional calls.

\begin{figure*}[t]
\centering\includegraphics[width=0.98\textwidth]{figures/operator.pdf}
\caption{(A) ID-based execution validates all $n$ judgments before emitting a ratio; direct plans skip the specification. (B) Illustrative complete, cached keyed outputs on one four-candidate pool. Output $a$ attains the exact reference ratio despite two candidate errors. ID 2 is routed to a pointwise judgment; disagreement does not reveal the shared candidate error at ID 3.}
\label{fig:operator}
\end{figure*}

We study these distinctions through an empirical audit of fixed-pool semantic aggregation. The target is the ratio under an explicit relevance mapping, not a latent notion of user satisfaction. An execution harness, retained under the artifact name IntentSat-2S, fixes candidate accounting and implements direct, specification-based, and selective plans. These plans are objects of comparison; specification extraction and third-on-disagreement checking are not presented as new algorithms. Elementary bounds explain what incomplete or inconsistent outputs permit us to infer, without certifying semantic correctness.

The study makes three contributions. First, a controlled audit crosses input-row order with output-key serialization while fixing candidate identities, chunk membership, prompt templates, and decoding hyperparameters; row order and the output-key grammar are the explicit interventions; repeated executions and same-template singletons provide diagnostic references. Second, comparisons on two product datasets and adapted biomedical and legal workloads expose where structural completion, candidate fidelity, and ratio loss diverge. Third, a trace-level cost audit compares disagreement checking with both exact equal-budget uniform checking and consensus-preserving AND/OR rules. The latter reveals that additional judgments can be unnecessary or harmful even when the routing signal looks useful. We report adverse results, distinguish new sampling from exposed-sample analyses, and preserve inputs, protocols, recorded responses, and verification scripts.

The practical question is therefore specific: for a fixed semantic aggregate, which execution checks provide evidence relevant to the target, and which merely add work? Our results answer this for the measured configurations; they do not establish a universally best predicate or a production optimizer.

\section{Background and Related Work}
\textbf{Semantic execution and batching.} LOTUS defines semantic operators and optimizes a reference algorithm under accuracy assumptions \cite{lotus}; Palimpzest explores cost--quality trade-offs \cite{palimpzest}, DocETL studies rewrites \cite{docetl}, and SemBench compares semantic engines \cite{sembench}. Sema already uses structured batching with fixed result cardinality and studies local scheduling \cite{sema}. NousDB calibrates batch sizes through sampled batch--singleton agreement and accounts for contextual interference \cite{nousdb}. Our fixed chunk limit is not a competing batch-size optimizer. We examine row order and output serialization jointly, then distinguish their effects from replay variation; our LOTUS comparisons cover a filter and one cascade, not its full optimizer.

\textbf{Repeated judgments and aggregate targets.} Run-to-run variance in semantic query processing, cancellation of candidate changes in COUNT, and a third vote only when two calls disagree have been studied directly \cite{runvariance}. Our selective rule is a separately prompted instance of this familiar binary majority pattern. We audit its value against two consensus-preserving rules and an exact equal-label-budget reference. Evergreen verifies semantic proportion claims \cite{evergreen}; SemCEB benchmarks model-dependent cardinalities \cite{semceb}. SemViews and SemReuse certify reuse against recorded or fixed oracles \cite{semviews,semreuse}. We evaluate human relevance proxies and do not challenge certificates conditional on such oracles.

\textbf{Incomplete outputs and contextual effects.} ThalamusDB and OLLA bound or progressively estimate partially evaluated semantic aggregates \cite{thalamus,olla}; query completeness and missing-data analysis provide broader foundations \cite{query_completeness,contingency}. Our operational policy abstains on unresolved labels, and the conditional intervals below serve as diagnostics. Semantic integrity constraints address grounding and soundness beyond output shape \cite{constraints}. Context sensitivity, permutation aggregation, and positional bias are also established \cite{lostmiddle,permcons,outputorder}; neither an ID parser nor our factorial audit identifies a model's latent reasoning. vLLM and SGLang can serve pointwise requests concurrently \cite{vllm,sglang}, so fewer logical calls do not by themselves imply lower latency.

\textbf{Relevance proxies.} Shopping Queries supplies human Exact, Substitute, Complement, and Irrelevant (E/S/C/I) labels \cite{esci}. We use E/S as positive and C/I as negative; WANDS uses a different Exact-only mapping \cite{wands}. These task definitions must accompany every result. Neither establishes purchase satisfaction, and results under one mapping cannot be silently transferred to the other.

\section{Target, Contract, and Diagnostic Quantities}
For query $q$, let $C_q=(c_1,\ldots,c_n)$ be a fixed candidate pool, $n>0$. Each $c_i$ is an indexed occurrence; repeated products remain separate positions. Under a specified reference predicate $s(q,c_i)\in\{0,1\}$, the target and a completed model estimate are
\begin{equation}
r_q=\frac{1}{n}\sum_i s(q,c_i),\qquad
\hat r_q=\frac{1}{n}\sum_i\hat s_i.
\end{equation}
The source-label mapping, pool construction, and text truncation define the task. This fixed-list ratio is not a catalog prevalence estimate or a measure of a different retrieval policy.

\textbf{Output contract.} Every occurrence must receive one valid binary label before a ratio is emitted. Array-form indexed JSON rejects missing, duplicate, or out-of-range IDs, duplicate keys, and extra fields. Keyed JSON requires the exact specified keys and Y/N values, then aligns by ID. Earlier plans restart IDs within chunks; the controlled audit preserves global pool-occurrence IDs in batches and singletons. Compact strings require exactly one Y/N character per input position and cannot detect swaps. Early wrappers trim whitespace; the controlled audit and biomedical/legal keyed runs validate raw responses. Retry limits and specification fallback are explicitly plan-specific; unresolved pools abstain.

\textbf{Two distinct losses.} For a complete pool, candidate error $e$ and ratio error satisfy
\begin{equation}
e=\frac{\mathrm{FP}+\mathrm{FN}}{n},\qquad
|\hat r_q-r_q|=\frac{|\mathrm{FP}-\mathrm{FN}|}{n}.
\label{eq:cancellation}
\end{equation}
Thus $e-|\hat r_q-r_q|=2\min(\mathrm{FP},\mathrm{FN})/n$. Cancellation is a known property of aggregate evaluation \cite{runvariance}, not an aggregate failure: an exact ratio achieves the ratio-only objective. Candidate fidelity is reported separately to diagnose the predicate and its suitability for subsequent filtering or provenance reuse. We do not substitute a candidate-level improvement for a ratio-level one.

\textbf{What partial output permits.} If $m$ accepted labels contain $p$ positives and at most $k$ errors, with integer $k$ and $0\leq k\leq m\leq n$, then
\begin{equation}
r_q\in\left[\frac{\max(0,p-k)}{n},\,
\frac{n-m+\min(m,p+k)}{n}\right].
\label{eq:bounds}
\end{equation}
Flipping accepted positives or negatives and assigning all missing labels to the corresponding extreme attains these endpoints over unrestricted occurrence labels. The error budget must come from outside the parser; without a justified smaller budget, $k=m$ gives $[0,1]$. These are conditional diagnostic bounds, not confidence intervals or emitted partial answers. The artifact also implements coupled-duplicate constraints and checks the combinatorics independently. Exchangeable calibration errors could support a marginal budget \cite{conformal}; our exposed, sequentially selected samples do not justify that guarantee.

\textbf{What disagreement permits.} For two complete aligned vectors $a,b$, let $D$ be their Hamming distance. At least one label is wrong at every disagreement, hence
\begin{equation}
e(a,s)+e(b,s)\geq D/n,\quad
\max\{e(a,s),e(b,s)\}\geq\lceil D/2\rceil/n.
\label{eq:stability}
\end{equation}
Agreement supplies no correctness certificate, and these inequalities do not identify the worse execution. The triangle inequality similarly lower-bounds the larger absolute ratio error by half the change in predicted ratio; unchanged counts can still accompany changed labels. These elementary relationships motivate separate reporting, rather than a new end-to-end accuracy guarantee.

Conditional on a model output, its aggregate is deterministic. Binary labels alone do not justify a binomial confidence interval: the pool is fixed and decisions can be dependent. We use query-cluster resampling to describe paired differences across the sampled pools and report completion on all attempts.

\section{Execution Plans}
Figure~\ref{fig:operator} summarizes the earlier plans; the controlled audit instead uses direct judgments, global IDs, and no retry. Query-only specification extraction returns a nonempty \texttt{product\_type} string and \texttt{required}/\texttt{preferred} string arrays. Despite instructions to use only query-supported constraints, it can invent requirements or omit a type. Validation checks shape, not truth.



Judgment prompts receive the query, products, and any specification. Indexed JSON links labels to IDs; compact codes follow input order. Invalid batches receive one retry. Invalid specifications cause abstention or direct compact fallback, by plan. Chunks contain at most eight candidates; all must validate before aggregation.

Selective execution combines two \emph{cached}, ID-aligned vectors $a,b$, retaining agreements and using a pointwise label $h_i$ at disagreements; unresolved labels cause abstention. Swapping these fixed vectors leaves their combination unchanged. Sorting an unchanged chunk by stable IDs makes its prompt presentation-invariant. Neither property guarantees accuracy, stability across different chunks/IDs, or deterministic recomputation.

Pointwise, direct-batch and two-stage execution require at least $n$, one and two calls, respectively; retries/fallback change costs. Call counts alone do not establish speed: batched prompts/JSON are longer, and pointwise calls can run concurrently. We record tokens and offline/local-service times.

\section{Experimental Method}
We organize the evaluation around three questions: \textbf{RQ1}, how validation, retries, and fallback affect completion and semantic quality; \textbf{RQ2}, how within-chunk row order and output-key serialization affect candidate labels and ratios beyond replay variation; and \textbf{RQ3}, whether additional disagreement judgments justify their cost against equal-budget and no-extra-judgment references. Later transfer and sensitivity checks probe the scope of the findings. Each follow-up is labeled by its design history; exposed-sample analyses are exploratory.

Table~\ref{tab:studies} separates study samples from repeated executions. Later models and audits reuse the listed pools; those repetitions do not add independent samples.

\begin{table}[t]
\caption{Study samples and their roles. Counts are attempted pools, before completion filtering. ACORD's two scenarios share 57 query IDs; the controlled WANDS audit reuses the 52 selective-execution pools. Selection rules and design history are detailed below.}
\label{tab:studies}
\centering\footnotesize
\begin{tabular}{lrl}
\toprule
Source / study & Pools & Main role \\
\midrule
ESCI derived pilot & 160 & Prompt development \\
ESCI derived validation & 500 & Completion and quality \\
ESCI original lists & 250 & Labeled-pool comparison \\
WANDS initial & 160 & Product-source transfer \\
WANDS holdout & 160 & Order audit \\
WANDS fallback & 96 & Cached-trace routing \\
WANDS selective & 52 & Selective execution \\
ESCI selective lists & 250 & Selective execution \\
ESCI controlled audit & 128 & Row/key-order factorial \\
TREC-COVID & 50 & Biomedical transfer \\
ACORD uniform / stress & 57 each & Legal transfer \\
\bottomrule
\end{tabular}
\end{table}

\subsection{Data Provenance and Pool Construction}
The public \texttt{nixiesearch/amazon-esci} conversion \cite{derived} uses US-English Shopping Queries small-version records. It maps E/S/C/I to $1/0.1/0.01/0$ and emits one row per E product with the query's non-E products, omitting other E products and original query/product IDs. Our test file matches the published SHA-256 at the recorded revision.

We retain the first eligible row per exact query string and reservoir-sample without replacement. Eligibility requires nonempty query/positive text, aligned negative texts/scores, and recognized codes. Each pool contains its E product and up to seven randomly selected non-E products, shuffled with a query-specific seed and truncated to 600 characters. This artificial composition can omit annotation evidence; string grouping can merge original query IDs. The estimand concerns constructed pools.

A 160-query pilot motivated compact output and specification fallback. A protocol frozen before the second run selected 500 other eligible test strings from the same source/model, establishing query-text disjointness rather than dataset independence. These pools contain @@NCAND@@ candidates (grades: @@GRADES@@). Reference labels map E/S to Y and C/I to N; they are relevance judgments, not purchase-satisfaction labels.

\subsection{Original Labeled Candidate Pools}
From hash-verified original Amazon Science tables at a recorded revision, we uniformly sample 250 US-English small-version test query IDs excluding earlier exact strings. All labeled examples successfully join products: 5,160 candidates (mean 20.64, range 8--79), with 2,293 E, 1,745 S, 242 C, and 880 I judgments. These same-source \emph{labeled lists} may omit unjudged/live results.

\subsection{Independent Wayfair Evaluation}
After observing ESCI, we add independently sourced Wayfair WANDS \cite{wands}. Its guidelines define Exact as a full match and Partial as an entity match missing a modifier; we map only Exact to Y. We collapse identical query--product labels, exclude 14 conflicting pairs and earlier ESCI query strings, and uniformly select 160 query IDs with at least 32 nonconflicting judgments. Uniformly sampling 32 judged products per query without grade stratification yields 5,120 candidates: 781 Exact, 3,282 Partial, and 1,057 Irrelevant. These labeled samples are not full retrieved pages.

Both models run four plans with unchanged contracts/limits and a full-match instruction. Truncated names, classes, features and descriptions omit images/pages available to annotators. The post-ESCI design, mapping and instruction confound attribution to dataset source.

\subsection{Disjoint WANDS Holdout}
@@HOLDOUT_METHOD@@

\subsection{Disagreement-Fallback Validation}
@@FALLBACK_METHOD@@

@@SELECTIVE_METHOD@@

@@ESCI_SELECTIVE_HOLDOUT_METHOD@@

\subsection{Controlled Row and Key Order Audit}
@@CONTROLLED_METHOD@@

\subsection{Post-hoc Local Service Test}
@@LOCAL_SERVICE_METHOD@@

\subsection{Plans and Measurements}
On identical shuffled pools, we compare (P) pointwise Y/N, (J) direct indexed JSON, (TJ) two-stage indexed JSON, (C) direct compact codes, and (TC) two-stage compact codes with direct fallback. Each invalid pointwise answer or batch gets one retry; TJ abstains on invalid specifications after retry, TC falls back. Temperature is zero; token caps are 4 for P, 192 for JSON, 160 for specifications, and 32 for compact codes. Original-pool runs use P/J/TJ/TC and chunks of at most eight. Pre-specified references are all-Y, all-N, and lexical Y iff every query word appears in the truncated text.

Candidate precision, recall, F1, and balanced accuracy use pooled confusion counts under each study's mapping; larger pools contribute more decisions. Balanced accuracy averages sensitivity and specificity. Ratio MAE averages $|\hat r_q-r_q|$ equally across complete queries. Completion uses every attempted pool; main paired quality comparisons use pools completed by all model plans. We separately label the two-JSON common-pool comparison. Derived-validation constant baselines are post hoc. Paired 95\% query-cluster bootstrap intervals use 4,000 resamples and describe this sampling procedure. They do not ensure population-level satisfaction; conditioning on completion can bias common subsets.

Primary inference uses Qwen3-4B-Instruct-2507 \cite{qwen}, one RTX 5090 (32 GB), bfloat16, vLLM 0.10.2, PyTorch 2.8.0/CUDA 12.8, a 4,096-token context, eager execution and at most 16 sequences. Offline times include chat-template rendering, generation, retries and fallback; they exclude loading and most other assembly/parsing. Logical prompt tokens do not measure cache recomputation; times are not online or hardware-independent latency. After Qwen, SmolLM2-1.7B-Instruct \cite{smollm} repeats the original 250-pool protocol with identical prompts, chunking, caps and GPU settings. Family, size, templates and tokenizer change together. The artifact retains scripts, outputs, model hashes and inputs.

@@THIRD_MODEL_METHOD@@

Our \emph{post-hoc supervised lexical baseline} is logistic regression on hashed query/product, shared-word and coverage features. It uses 50,000 ESCI train examples from 18,490 strings, excluding all case-folded, trimmed evaluation queries; 357 product IDs still overlap the 250-pool test sample. Training uses E/S versus C/I; evaluation uses the same truncated text. Features and threshold were fixed before the first run; its iteration-limit warning led to a higher limit, with both outputs retained. Supervision and shared dataset provenance distinguish it from the zero-shot plans.

@@WANDS_SUPERVISED_METHOD@@

@@WANDS_EMBEDDING_METHOD@@

Post hoc, LOTUS 1.2.4 runs its actual \texttt{sem\_filter} with the same Qwen model, data and relevance proxy over a local vLLM API. Prompts, parsing, caps and online execution differ. We retain answers/decisions and count its configured True defaults on unparseable answers. One filter does not represent optimized LOTUS.

@@LOTUS_CASCADE_METHOD@@

Post hoc, we remove the four-item compact example, inspect strict/delimiter-tolerant parsing, and compare the revised prompt with/without V1 guided decoding to the actual chunk length. Saved specifications/fallbacks, candidates and order stay fixed. An earlier V0 attempt did not visibly enforce the constraint and is retained. We then reverse within chunks, reassign position IDs, and realign labels; subsequent identical replay and stable-ID reversal probe that confound. These fixed-sequence audits reuse pools.

After inspecting the holdout, we reuse its saved specifications/pools for an exploratory keyed-JSON audit. Guided decoding fixes ascending local-ID keys with Y/N values. Original, replay and stable-ID-reversed responses are parsed by ID. Prompt, grammar and token cap differ from compact output; this is not another confirmatory experiment.

\textbf{Matched-label-budget ablation.} Post hoc, on the 52 WANDS and 250 ESCI selective pools for both models, we fix all keyed and full-pointwise traces. For each pool, both virtual plans replace exactly $k$ canonical labels, where $k$ is its observed disagreement count: either the disagreeing positions or a uniform $k$-subset. We compute the latter's expected MAE exactly from hypergeometric counts. Paired query-bootstrap intervals use 4,000 draws. This conditions on observed budgets and saved outputs; tokens, scheduling and fresh judgments can differ. We also combine the same two keyed vectors by AND (N on disagreement) and OR (Y on disagreement). Both preserve agreements and avoid every additional pointwise call. Both rules are reported on every sample, without selecting a rule using test labels. Reference-defined class-composition strata are post-hoc diagnostics.

\section{Results}
\subsection{RQ1: Completion and Semantic Quality}
\textbf{Development and validation.} The 160-query development pilot motivated compact output and specification fallback. A separate 500-query validation jointly completes 424 pools across five plans. On that common subset, the two-stage-JSON minus pointwise F1 interval is $[-0.250,-0.195]$; all-Y minus pointwise MAE spans zero, $[-0.017,0.051]$. The positive-heavy construction and conditioning on completion limit interpretation. Of 1,668 substitute examples, pointwise accepts 1,169 versus 548 for two-stage JSON; of 236 complements, the plans reject 79 versus 160. A stricter specification can reject both unwanted complements and valid substitutes under the E/S mapping.

On all 500 attempts, pointwise makes 3,582 calls in 27.25 seconds, whereas direct indexed JSON makes 501 calls in 57.23 seconds. Fewer calls are slower in this offline batch. Two-stage JSON has 31 invalid specifications and two other failed pools; compact fallback handles those specifications but still fails 30 pools. Different prompts and formats prevent isolating a fallback effect. The artifact retains all attempted outputs, timings, and per-plan metrics.

\subsection{Original Labeled Candidate Pools and Comparators}
@@FULL_RESULTS@@

\begin{table*}[t]
\caption{Original labeled query pools. Quality is computed on the @@FULL_COMMON@@ pools completed by every Qwen plan; completion, calls, token counts, and time cover all 250 attempted pools. The constant and exact lexical baselines are pre-specified; LOTUS and the trained classifier were added post hoc. LOTUS time is excluded because it uses an online API.}
\label{tab:full}
\centering\small
\begin{tabular}{lrrrrrrrr}
\toprule
Plan & Complete & F1 & Bal. acc. & Ratio MAE & Calls & Input tok. & Output tok. & Time (s) \\
\midrule
@@FULL_TABLE_ROWS@@
\bottomrule
\end{tabular}
\end{table*}

@@FULL_INTERPRETATION@@

\textbf{Truncation sensitivity.} @@ESCI_TRUNCATION_RESULTS@@

\begin{figure*}[t]
\centering\includegraphics[width=0.98\textwidth]{figures/quality.pdf}
\caption{Points show pooled F1 and query-macro ratio MAE on pools completed by all four Qwen plans; whiskers are 95\% query-bootstrap intervals. LOTUS and the trained comparator are post hoc; vertical lines are pre-specified baselines. Pools do not measure user satisfaction.}
\label{fig:quality}
\end{figure*}












\textbf{LOTUS filter.}
LOTUS completes all 250 ESCI pools, with no parser-default decisions among 5,160 candidates. On the common 207, its F1/MAE is 0.556/0.440; LOTUS minus pointwise intervals are $[-0.301,-0.217]$ for F1 and $[0.192,0.276]$ for MAE. Different prompts, contracts and online execution prevent attribution to the engine alone.

\textbf{Supervised lexical reference.}
The trained lexical reference has F1/MAE 0.836/0.182 on the same 207 pools. Its differences from pointwise have intervals spanning zero. Supervision, shared source provenance and 357 overlapping product IDs distinguish it from the zero-shot plans; the original nonconverged fit is retained alongside the converged run.

\textbf{Model sensitivity.}
SmolLM2 completes 250/250 pointwise pools but labels all 5,160 candidates N; F1 is 0 and ratio MAE 0.785. Every compact-fallback pool has an unresolved chunk, leaving the pre-specified four-plan, two-model intersection empty. Its post-hoc two-JSON comparison conditions on completion and does not repair this failure. Untuned prompts, caps, model family, tokenizer and chat template change together, so the result diagnoses this configuration rather than an isolated architecture effect.

\subsection{Wayfair Results and Format Audits}
@@WANDS_RESULTS@@

\begin{table*}[t]
\caption{Qwen on sampled WANDS pools using Exact as the positive label. Quality is conditional on the @@WANDS_COMMON@@ pools completed by all four model plans; completion, calls, token counts, and offline time cover all 160 attempted pools. Constant and lexical comparators are deterministic.}
\label{tab:wands}
\centering\small
\begin{tabular}{lrrrrrrrr}
\toprule
Plan & Complete & F1 & Bal. acc. & Ratio MAE & Calls & Input tok. & Output tok. & Time (s) \\
\midrule
@@WANDS_TABLE_ROWS@@
\bottomrule
\end{tabular}
\end{table*}


SmolLM2's P/J/TJ/TC plans complete 160/148/130/0 WANDS pools. Pointwise emits N for all 5,120 candidates (F1 0; MAE 0.153). The four-plan common subset and the pre-specified two-model intersection are empty. Completion therefore permits degenerate predictions; format failure can also remove every pool from a paired comparison.

\begin{figure*}[t]
\centering\includegraphics[width=0.98\textwidth]{figures/model_sensitivity.pdf}
\caption{Post-hoc model sensitivity: selective minus keyed ratio MAE with unadjusted 95\% query-cluster bootstrap intervals. Annotations show keyed $\to$ selective MAE; panel titles show common-complete pool counts. In this figure, all 7B product specifications failed strict validation and used direct fallback; TREC uses official narratives. Models, tokenizers and prompts differ, so the contrasts do not isolate model size.}
\label{fig:model_sensitivity}
\end{figure*}

\textbf{Format intervention.} SmolLM2's compact prompt contains a four-item example for chunks of up to eight. Removing it still leaves 0/250 ESCI and 0/160 WANDS complete; tolerant parsing recovers only 22 ESCI pools (F1 0.015). Guided decoding subsequently completes all pools. Its F1 remains 0.253/0.092 and MAE 0.643/0.149 on ESCI/WANDS, compared with 0.215 for ESCI all-Y and 0.153 for WANDS all-N. For Qwen, guided-minus-unguided F1 intervals span zero on common pools ($[-0.006,0.002]$ and $[-0.005,0.009]$). These exposed-sample interventions change format availability without establishing semantic accuracy.

\subsection{RQ2: Row Order, Output Keys, and Replay}
\textbf{Motivating audits.} On the pre-specified 160-pool WANDS holdout, stable-ID within-chunk reversal changes 1,733/5,120 Qwen compact labels (33.8\%; 95\% query interval $[31.8,35.9]\%$), while identical replay changes none. SmolLM2 changes 643/5,120 (12.6\%). Both models complete all conditions. Qwen's mean absolute ratio change is 0.048; its original F1/MAE is 0.388/0.177, versus 0.595/0.115 for the all-query-words reference. Three later permutations change 34.6--35.1\% of Qwen and 8.2--8.7\% of SmolLM2 labels; their replay controls change none.

Keyed JSON on those exposed pools changes 652/5,120 Qwen labels (12.7\%; $[11.3,14.3]\%$) and 596/5,120 SmolLM2 labels (11.6\%; $[11.0,12.3]\%$) under stable-ID reversal, with zero replay changes. Qwen's original keyed F1/MAE is 0.574/0.148. Compact and keyed runs differ in prompts, grammar, and caps, so their gap does not isolate serialization. The compact instruction requires input-order labels; reinterpreting them in numeric-ID order would change disagreement rates but would contradict that instruction. Earlier ID-reassignment controls and raw traces remain in the artifact.

Replay variation also depends on runtime: a later isolated, local-copy replay on 16 exposed pools changes 13/512 labels between identical prompts under changed runtime/batch settings. This motivates interleaved repeated conditions in the controlled audit, instead of treating a previous zero replay rate as a universal property.

\textbf{Controlled factorial audit.} @@CONTROLLED_RESULTS@@

\subsection{RQ3: Disagreement Checking and Cost}
On the earlier 96-pool WANDS sample, all conditions complete; cached keyed/pointwise-fallback MAEs are 0.124/0.106 (paired 95\% query interval for the difference $[-0.0335,-0.0029]$), versus 0.094 for lexical fallback. On the subsequent 52 pools, actual selective execution changes Qwen's MAE from 0.136 to 0.120 ($[-0.0325,-0.0018]$), using 195 pointwise checks and 663 calls overall. A query--product reweighting sensitivity interval $[-0.0395,0.0064]$ includes zero and changes the estimand. Four of the 195 checker outputs differ in a later full-pointwise run. SmolLM2 changes MAE from 0.190 to 0.195 ($[-0.0012,0.0120]$). Both models complete all 52 pools; Table~\ref{tab:fallback} includes the task-specific references.

On the 250 ESCI selective pools, both models complete every condition. Qwen's keyed/selective/pointwise MAEs are 0.407/0.373/0.299; selective minus keyed has query interval $[-0.045,-0.023]$. All-Y and the supervised reference reach 0.185/0.183. Selective uses 2,133 versus 4,690 full-pointwise calls, but 1.52M versus 1.05M logical input tokens. SmolLM2's keyed/selective MAEs are 0.802/0.813 ($[0.007,0.014]$). Excluding the 18 text-overlap queries preserves the rounded differences, $-0.034/+0.011$. Rescoring fixed Qwen outputs with only E positive changes its difference to $+0.008$ ($[-0.004,0.019]$); that post-hoc target differs from the E/S prompt. Calls and tokens do not establish a speedup.

\textbf{Ratio-error cancellation.} Post hoc, Qwen's selective candidate error/ratio MAE is 0.150/0.120 on WANDS and 0.420/0.373 on the ESCI selective sample. Seven of ten ESCI pools with exact selective ratios contain candidate errors; these ten still succeed at the ratio-only objective. All-Y has equal candidate error and MAE, 0.185, because its errors have one direction. Section III's identity explains these query-macro differences.

@@ROUTING_BUDGET_RESULTS@@

\textbf{Consensus-preserving references.} The comparisons below use the actual sparse-checker outputs. The uniform-budget analysis instead uses saved full-pointwise outputs: on WANDS/Qwen these give routed MAE 0.118, versus 0.120 for actual selective execution. On ESCI/Qwen the two checker traces agree. 
Post hoc, we retain the same two keyed outputs and compare AND (N on disagreement)
and OR (Y on disagreement), reporting both policies on every sample
(Table~\ref{tab:consensus}). On the 250-pool ESCI selective sample, Qwen's OR has MAE 0.338 and F1 0.704,
versus 0.373 and 0.677 for selective checking. The selective-minus-OR MAE
difference is 0.0354, with an unadjusted 95\% paired query-bootstrap interval
$[0.0273,0.0438]$. OR uses 1,394 rather than 2,133 trace-derived calls and
1.345M rather than 1.517M logical input tokens; its balanced accuracy is
slightly lower, 0.641 versus 0.649. On WANDS, AND has lower MAE, but the
selective-minus-AND interval $[-0.0174,0.0294]$ includes zero and selective
has higher balanced accuracy. SmolLM2's selective vectors equal AND on both
samples because all disagreement checks return N. These virtual comparisons
show that an advantage over uniform checking need not justify extra judgments;
they do not measure latency or identify a universally best tie-breaker.

\begin{table}[t]
\caption{Post-hoc consensus references on the same complete saved traces.
Entries are query-macro ratio MAE. AND and OR use the same two keyed passes
and specification as selective checking; the final column counts pointwise
calls they avoid. Both policies are reported on all samples; neither is chosen as a policy using test labels.}
\label{tab:consensus}
\centering\footnotesize
\begin{tabular}{lrrrr}
\toprule
Sample / model & AND & OR & Selective & Calls avoided \\
\midrule
WANDS / Qwen & 0.115 & 0.173 & 0.120 & 195 \\
WANDS / SmolLM2 & 0.195 & 0.194 & 0.195 & 192 \\
ESCI / Qwen & 0.477 & 0.338 & 0.373 & 739 \\
ESCI / SmolLM2 & 0.813 & 0.770 & 0.813 & 190 \\
\bottomrule
\end{tabular}
\end{table}

\textbf{Consensus-preserving count diagnostic.}
Post hoc, let $L$ and $U$ be the positive counts under AND and OR of the two keyed
outputs, and $G$ the reference count. Every rule retaining their agreements has
its positive count in $[L,U]$; the minimum ratio error over all such rules is
$\max\{L-G,0,G-U\}/n$. On ESCI, $G\geq U$ in 221/250 Qwen pools;
the query-macro floor is 0.322, with OR and actual selective excess losses
of 0.016 and 0.051. All 250 SmolLM2 pools satisfy $G\geq U$; its floor is
0.770 and selective excess is 0.043. On WANDS, Qwen has $G\leq L$ in 35/52
pools, with floor 0.085 and AND/selective excess 0.029/0.034. SmolLM2 has
20 lower-side, 8 interior and 24 upper-side pools; its floor is 0.137 and
selective excess is 0.058. Collapsed exact cases $L=U=G$ are classified separately;
none occur. Independently checked per-pool records retain all four settings.
These reference-dependent diagnostics expose limits of retaining agreement;
they neither select an online endpoint nor bound plans that reconsider agreed labels.

\textbf{Composition of the budget contrast.}
Post hoc, Qwen's cached routed-minus-expected-uniform MAE on mixed-label pools
is $+0.0005$ on WANDS ($n=32$, unadjusted 95\% query interval
$[-0.0122,0.0131]$) and $-0.0043$ on ESCI ($n=154$,
$[-0.0135,0.0045]$). The 20 all-negative WANDS pools contribute $-0.0120$
to the overall $-0.0117$ contrast; the 96 all-positive ESCI pools contribute
$-0.0087$ to $-0.0114$. These reference-defined subgroups describe the exposed
sample; they neither supply an online routing rule nor establish no effect
or a significant interaction on mixed pools.

\begin{figure*}[t]
\centering\includegraphics[width=0.98\textwidth]{figures/routing_budget.pdf}
\caption{Post-hoc cached-trace routing minus the exact expected error of uniform checking, with equal per-pool pointwise-label budgets derived from keyed disagreements. Both use the same saved full-pointwise judgments; the reference averages all subsets of that size. Whiskers are unadjusted 95\% query-bootstrap intervals. Candidate error is query-macro. This does not compare newly executed pipelines or total resources.}
\label{fig:routing_budget}
\end{figure*}

\textbf{Two-order references.} On 52 WANDS pools, Qwen's mean of two keyed ratios has MAE 0.139 versus 0.120 selective (paired 95\% selective-minus-midpoint interval $[-0.033,-0.007]$); the midpoint gives no new candidate labels. A virtual tie-breaker applies all-query-words labels on keyed disagreements without extra LLM calls: MAE/F1 is 0.103/0.657 versus 0.120/0.639 selective (selective-minus-lexical MAE interval $[-0.002,0.038]$); selective adds 195 calls. On ESCI, Qwen's lexical/selective MAE is 0.459/0.373 and F1 is 0.588/0.677 (selective-minus-lexical MAE interval $[-0.100,-0.072]$), with 739 extra calls. These unadjusted post-hoc comparisons do not measure online or lexical CPU time.

@@LOTUS_CASCADE_RESULTS@@

\begin{table*}[t]
\caption{Disagreement fallback and references on two disjoint WANDS samples.}
\label{tab:fallback}
\centering\small
\begin{tabular}{lrrrrrr}
\toprule
Plan & Complete & F1 & Bal. acc. & Ratio MAE & Gen. calls & Input tokens \\
\midrule
@@FALLBACK_TABLE_ROWS@@
\bottomrule
\end{tabular}
\par\smallskip
\begin{minipage}{0.94\textwidth}\footnotesize
Quality uses common-complete pools; samples are not directly comparable. The 96-query fallbacks use saved traces; the 52-query scheduler sends actual requests. Dashes omit encoder inference, not zero cost; reference training is excluded. CV uses training folds; LOTUS totals include helper, calibration and routed main calls. Token totals combine logical lengths across models, not latency.
\end{minipage}
\end{table*}

\subsection{Local Service Results}
@@LOCAL_SERVICE_RESULTS@@

\begin{table}[t]
\caption{Local HTTP workload on the existing 52 WANDS queries. Each range spans two repetitions; completion lists both counts out of 52. Latency uses only complete queries and excludes client admission, startup and warmup.}
\label{tab:local_service}
\centering\footnotesize
\begin{tabular}{lrrr}
\toprule
Plan (active queries) & p50 (s) & p95 (s) & Complete \\
\midrule
@@LOCAL_SERVICE_ROWS@@
\bottomrule
\end{tabular}
\end{table}

\textbf{Post-hoc WANDS references.} On exposed 160/96/52 WANDS queries, ratio-tuned lexical/embedding MAEs are 0.172/0.128/0.189 and 0.180/0.122/0.181. Their 52-query differences from Qwen have unadjusted intervals $[0.013,0.130]$ and $[0.005,0.123]$; embedding product reweighting gives $[-0.004,0.126]$ on a changed estimand. F1-tuned results are in Table~\ref{tab:fallback}.

\subsection{Third-Model Sensitivity}
In unconstrained 7B runs, all 250/52 product specifications fail validation and use direct fallback. ESCI keyed/selective MAEs are 0.674/0.622 versus 0.185 all-Y; WANDS MAEs are 0.170/0.163. On TREC, MAE rises from 0.100 to 0.128 (paired interval $[0.018,0.039]$). Figure~\ref{fig:model_sensitivity} reports unadjusted, descriptive comparisons on these exposed samples.

Schema guidance removes specification fallback on both product samples (0/250 ESCI, 0/52 WANDS). Selective MAE changes from 0.622 to 0.648 on ESCI (unadjusted paired interval $[0.012,0.040]$) and from 0.163 to 0.150 on WANDS ($[-0.028,-0.001]$). Product-reweighted ranges $[0.008,0.045]$ and $[-0.033,0.004]$ change the estimand. Valid format does not certify meaning.

\subsection{Conditional Bounds and Calibration}
@@BOUNDS_DIAGNOSTIC@@

\subsection{Cross-Domain Transfer Check}
@@TREC_TRANSFER@@


\textbf{Legal transfer.} After earlier outcomes, an ACORD protocol \cite{acord} was fixed before sampling/inference for 57 queries, each with 32-clause uniform and label-constructed stress pools. Scores 3/4 map to Y; some zeros are source-expanded. Queries serve as specifications. Length inspection set 2,400-character prefixes before inference; 195/1,824 uniform and 251/1,824 stress occurrences remain truncated. Qwen completes 114/114 pools. Stress keyed/selective MAEs are 0.168/0.137 (difference -0.0307; unadjusted query-bootstrap $[-0.0406,-0.0214]$); pointwise/lexical MAEs are 0.099/0.123. Selective uses 1,234 calls versus 3,648 pointwise across scenarios. A post-hoc stress-pool four-minus-eight-item recovery comparison has MAE-difference interval $[-0.0164,0.0060]$, confounded by partitions, grammars and conflicting dependencies. Uniform selective MAE 0.051 exceeds all-N 0.008 (15 positives); stress has 359 positives. Untuned SmolLM2, following Qwen on these pools, changes 1/3,648 labels under reversal and returns all-N stress selective outputs (MAE 0.197; F1 0). Pretraining exposure, reuse, label-dependent sampling and truncation limit generalization.

\section{Limitations and Research Implications}
\textbf{Biased candidate pools.} These are judged lists, not live results. Original/selective ESCI have 172/250 and 197/250 size-16 pools; the latter has text overlaps and 96 all-positive pools, separately from the 128-pool controlled audit. WANDS follow-ups have 42/96 and 20/52 all-negative pools. Reference-defined groups cannot route live queries.

\textbf{Semantic and data limitations.} Truncation, missing attributes, ambiguous queries and labels limit conclusions; relevance is not purchase satisfaction. Pretraining exposure is unknown; query bootstrap cannot repair sampling bias.

\textbf{Systems scope.} Four datasets, three models and one GPU cannot establish deployment performance. Earlier reversal audits confound presentation choices; the controlled factorial fixes several of them but does not reveal latent reasoning or isolate hardware effects. Production traffic, unjudged candidates and fresh-machine reproduction remain untested.

\textbf{Design implications.} Validate fixed-pool accounting before aggregation. Report candidate and ratio losses separately, retain constant baselines, and measure the full pipeline. Compare equal-budget allocations and AND/OR rules before adding a checker. Disagreement need not identify the most useful positions; agreement certifies no correctness. Error cancellation can validly succeed at a ratio-only task.

\section{Conclusion}
Semantic aggregate evaluation needs fixed candidate accounting, explicit target losses and references with stated budgets. Completion cannot certify accuracy. Disagreement checking can beat a keyed pass or uniform allocation yet trail OR; an all-negative checker reduces to AND. Controlled contrasts compare row/key-order variation with replay. Evaluate the full plan against inexpensive references, with separate evidence for any claimed semantic or latency benefit.

\section*{Acknowledgments and AI-Generated Content Disclosure}
OpenAI Codex (GPT-6) substantially drafted and revised Sections I--VIII, wrote evaluation and analysis scripts and the code for all manuscript figures, and assisted literature review and interpretation. Qwen3-4B-Instruct-2507, Qwen2.5-7B-Instruct and SmolLM2-1.7B-Instruct generated Section VI judgments; Qwen3-4B also generated the artifact's synthetic format-check outputs. Qwen3-Embedding-0.6B generated vectors under Section V's protocol. The authors are responsible for the correctness and originality of the text, code, figures, sources and results.

\IEEEtriggeratref{15}
\begin{thebibliography}{99}
\bibitem{lotus} L. Patel, S. Jha, M. Pan, H. Gupta, P. Asawa, C. Guestrin, and M. Zaharia, ``Semantic operators and their optimization: Enabling LLM-based data processing with accuracy guarantees in LOTUS,'' \emph{Proc. VLDB Endow.}, vol. 18, no. 11, pp. 4171--4184, 2025, doi: 10.14778/3749646.3749685.
\bibitem{docetl} S. Shankar, T. Chambers, T. Shah, A. G. Parameswaran, and E. Wu, ``DocETL: Agentic query rewriting and evaluation for complex document processing,'' \emph{Proc. VLDB Endow.}, vol. 18, no. 9, pp. 3035--3048, 2025, doi: 10.14778/3746405.3746426.
\bibitem{palimpzest} C. Liu \emph{et al.}, ``Palimpzest: Optimizing AI-powered analytics with declarative query processing,'' in \emph{Proc. CIDR}, 2025.
\bibitem{sembench} J. Lao \emph{et al.}, ``SemBench: A benchmark for semantic query processing engines,'' \emph{Proc. VLDB Endow.}, vol. 19, no. 8, pp. 1754--1767, 2026, doi: 10.14778/3811243.3811249.
\bibitem{sema} K. Qi \emph{et al.}, ``Sema: A high-performance system for LLM-based semantic query processing,'' arXiv:2603.11622, 2026.
\bibitem{nousdb} C. Tsapelas, G. Koutrika, and Y. Ioannidis, ``NousDB: A fast query engine for semantic analytics,'' arXiv:2610.28551, 2026.
\bibitem{evergreen} A. W. Lee, B. Han, S. Sen, S. Yeom, U. Cetintemel, and A. Datta, ``Evergreen: Efficient claim verification for semantic aggregates,'' arXiv:2604.26180, 2026.
\bibitem{semceb} A. Zimmerer, C. K\"uhn, Y. Li, M. Stoian, R. Borovica-Gajic, and A. Kipf, ``SemCEB: A cardinality estimation benchmark for semantic operators,'' arXiv:2606.23081, 2026.
\bibitem{semviews} O. Hassanzadeh and D. Subramanian, ``Semantic views: Certified reuse of LLM labels across semantic queries,'' arXiv:2610.12768, 2026.
\bibitem{semreuse} P. T. Tran-Truong, T. Kieu, X.-B. Le, S. X. Ha, D. Kieu, B. Yang, and C. S. Jensen, ``SemReuse: Certified reuse of cached natural-language views for semantic operators,'' arXiv:2610.16035, 2026.
\bibitem{thalamus} S. Jo and I. Trummer, ``ThalamusDB: Approximate query processing on multi-modal data,'' \emph{Proc. ACM Manag. Data}, vol. 2, no. 3, Art. 186, 2024, doi: 10.1145/3654989.
\bibitem{olla} C. Hui, W. Lu, Y. Gao, L. Xiong, Y. Wang, and Y. Chen, ``LLM-driven online aggregation for unstructured text analytics,'' in \emph{Proc. DASFAA}, 2026, pp. 578--594, doi: 10.1007/978-981-92-0366-6\_35.
\bibitem{query_completeness} S. Razniewski and W. Nutt, ``Completeness of queries over incomplete databases,'' \emph{Proc. VLDB Endow.}, vol. 4, no. 11, pp. 749--760, 2011, doi: 10.14778/3402707.3402715.
\bibitem{contingency} X. Liang, Z. Shang, S. Krishnan, A. J. Elmore, and M. J. Franklin, ``Fast and reliable missing data contingency analysis with predicate-constraints,'' in \emph{Proc. SIGMOD}, 2020, pp. 285--295, doi: 10.1145/3318464.3389785.
\bibitem{constraints} A. W. Lee \emph{et al.}, ``Semantic integrity constraints: Declarative guardrails for AI-augmented data processing systems,'' \emph{Proc. VLDB Endow.}, vol. 18, no. 11, pp. 4073--4080, 2025, doi: 10.14778/3749646.3749677.
\bibitem{lostmiddle} N. F. Liu \emph{et al.}, ``Lost in the middle: How language models use long contexts,'' \emph{Trans. Assoc. Comput. Linguistics}, vol. 12, pp. 157--173, 2024, doi: 10.1162/tacl\_a\_00638.
\bibitem{permcons} R. Tang, C. Zhang, X. Ma, J. Lin, and F. Ture, ``Found in the middle: Permutation self-consistency improves listwise ranking in large language models,'' in \emph{Proc. NAACL}, 2024, pp. 2327--2340, doi: 10.18653/v1/2024.naacl-long.129.
\bibitem{outputorder} P. Schilcher \emph{et al.}, ``Characterizing positional bias in large language models: A multi-model evaluation of prompt order effects,'' in \emph{Findings of EMNLP}, 2025, pp. 20643--20664, doi: 10.18653/v1/2025.findings-emnlp.1124.
\bibitem{runvariance} O. Hassanzadeh and D. Subramanian, ``Same query, different answer: Run-to-run variance in semantic query processing,'' arXiv:2610.13823, 2026.
\bibitem{vllm} W. Kwon \emph{et al.}, ``Efficient memory management for large language model serving with PagedAttention,'' in \emph{Proc. SOSP}, 2023, pp. 611--626, doi: 10.1145/3600006.3613165.
\bibitem{sglang} L. Zheng \emph{et al.}, ``SGLang: Efficient execution of structured language model programs,'' \emph{Adv. Neural Inf. Process. Syst.}, vol. 37, pp. 62557--62583, 2024, doi: 10.52202/079017-2000.
\bibitem{esci} C. K. Reddy \emph{et al.}, ``Shopping Queries Dataset: A large-scale ESCI benchmark for improving product search,'' arXiv:2206.06588, 2022.
\bibitem{derived} Nixiesearch, ``Amazon ESCI dataset in nixietune format,'' revision c092c898, 2024. [Online]. Available: \url{https://huggingface.co/datasets/nixiesearch/amazon-esci}
\bibitem{conformal} R. F. Barber, E. J. Cand\`es, A. Ramdas, and R. J. Tibshirani, ``Conformal prediction beyond exchangeability,'' \emph{Ann. Statist.}, vol. 51, no. 2, pp. 816--845, 2023, doi: 10.1214/23-AOS2276.
\bibitem{wands} Y. Chen, S. Liu, Z. Liu, W. Sun, L. Baltrunas, and B. Schroeder, ``WANDS: Dataset for product search relevance assessment,'' in \emph{Proc. ECIR}, 2022, pp. 128--141, doi: 10.1007/978-3-030-99736-6\_9.
\bibitem{pigeonhole} A. B. Owen, ``The pigeonhole bootstrap,'' \emph{Ann. Appl. Stat.}, vol. 1, no. 2, pp. 386--411, 2007, doi: 10.1214/07-AOAS122.
\bibitem{qwen} Qwen, ``Qwen3-4B-Instruct-2507,'' model card, 2025. [Online]. Available: \url{https://huggingface.co/Qwen/Qwen3-4B-Instruct-2507}
\bibitem{smollm} Hugging Face, ``SmolLM2-1.7B-Instruct,'' model card, 2024. [Online]. Available: \url{https://huggingface.co/HuggingFaceTB/SmolLM2-1.7B-Instruct}
\bibitem{qwen25} Qwen, ``Qwen2.5-7B-Instruct,'' model card, 2024. [Online]. Available: \url{https://huggingface.co/Qwen/Qwen2.5-7B-Instruct}
\bibitem{qwen_embed} Qwen, ``Qwen3-Embedding-0.6B,'' model card, 2025. [Online]. Available: \url{https://huggingface.co/Qwen/Qwen3-Embedding-0.6B}
\bibitem{treccovid} K. Roberts \emph{et al.}, ``Searching for scientific evidence in a pandemic: An overview of TREC-COVID,'' \emph{J. Biomed. Inform.}, vol. 121, p. 103865, 2021, doi: 10.1016/j.jbi.2021.103865.
\bibitem{acord} S. H. Wang \emph{et al.}, ``ACORD: An expert-annotated retrieval dataset for legal contract drafting,'' in \emph{Proc. ACL}, 2025, pp. 24739--24762, doi: 10.18653/v1/2025.acl-long.1206.
\end{thebibliography}
\end{document}
